%%%%%%%%%%%%%%%%%%%%%%%%%%%%%%%%%%%%%%%%%%%%%%%%%%%%%%%%%%%%%%%%%%%%%%%%%%%%%%%%%%
\begin{frame}[fragile]\frametitle{}
\begin{center}
{\Large राकेश देशपांडे --- ७-दिवसीय ध्यान अभ्यासक्रम}
\end{center}

{\tiny (Ref: राकेश देशपांडे --- ध्यान अभ्यासक्रम, YouTube)}

\end{frame}

%%%%%%%%%%%%%%%%%%%%%%%%%%%%%%%%%%%%%%%%%%%%%%%%%%%%%%%%%%%%%%%%%%%%%%%%%%%%%%%%%%
%% दिवस १: ध्यानाविषयी मिथके
%%%%%%%%%%%%%%%%%%%%%%%%%%%%%%%%%%%%%%%%%%%%%%%%%%%%%%%%%%%%%%%%%%%%%%%%%%%%%%%%%%

%%%%%%%%%%%%%%%%%%%%%%%%%%%%%%%%%%%%%%%%%%%%%%%%%%%%%%%%%%%
\begin{frame}[fragile]\frametitle{दिवस १ --- मिथक १: विचार थांबतात}

	\begin{itemize}
	\item ध्यानात विचार थांबवायला लागतात --- हे चुकीचे आहे
	\item मन सतत विचार निर्माण करत राहते --- ध्यानातही हे थांबत नाही
	\item ध्यान म्हणजे विचारांना जबरदस्तीने दडपणे नव्हे
	\item बियाणे पेरल्यावर, पाणी घातल्यावर, झाड आपोआप वाढते
	\item त्याचप्रमाणे योग्य वातावरणात मन आपोआप निवळते --- विचार आपोआप शांत होतात
	\end{itemize}

{\tiny (Ref: राकेश देशपांडे --- दिवस १)}

\end{frame}

%%%%%%%%%%%%%%%%%%%%%%%%%%%%%%%%%%%%%%%%%%%%%%%%%%%%%%%%%%%
\begin{frame}[fragile]\frametitle{दिवस १ --- मिथक २ व ३}

	\begin{itemize}
	\item \textbf{मिथक २}: निर्विचार अवस्था = शांती --- हेही खरे नाही
	\item शांतता म्हणजे केवळ विचार थांबणे नाही
	\item मानसिक संतुलन, जागरूकता आणि जीवनशैलीवर शांती अवलंबून
	\item ओव्हर थिंकिंगमुळे मन अस्थिर होते --- ते कमी होणे म्हणजे शांती
	\item \textbf{मिथक ३}: गाइडेड मेडिटेशन / म्युझिक मेडिटेशन = ध्यान --- नव्हे
	\item ध्यान म्हणजे अंतर्मुख होणे --- बाह्य आदेशावर प्रतिक्रिया देणे नव्हे
	\item योगनिद्रा = प्रत्याहार (ध्यानाचा पहिला टप्पा), प्रत्यक्ष ध्यान नव्हे
	\end{itemize}

{\tiny (Ref: राकेश देशपांडे --- दिवस १)}

\end{frame}

%%%%%%%%%%%%%%%%%%%%%%%%%%%%%%%%%%%%%%%%%%%%%%%%%%%%%%%%%%%
\begin{frame}[fragile]\frametitle{दिवस १ --- मिथक ४ + साधना}

	\begin{itemize}
	\item \textbf{मिथक ४}: चालत/धावत असताना ध्यान = अनुपस्थित मन (Absent-mindedness)
	\item मन विचारांमध्ये गुरफटलेले असते --- जागरूकतेचा अभाव असतो
	\item ध्यानासाठी निश्चित वेळ आणि स्थिर जागा आवश्यक
	\item झोप आली तर शरीराला पुरेशी विश्रांती मिळालेली नाही --- ७--८ तास झोप आवश्यक
	\end{itemize}

	\vspace{2mm}
	\textbf{साधना:} २० मिनिटे --- सर्व इंद्रियांनी आजूबाजूचे निरीक्षण करा \\
	(स्पर्श, चव, दृष्टी, वास, ऐकणे --- प्रत्येकी ५ मिनिटे; डोळे उघडे)

{\tiny (Ref: राकेश देशपांडे --- दिवस १)}

\end{frame}

%%%%%%%%%%%%%%%%%%%%%%%%%%%%%%%%%%%%%%%%%%%%%%%%%%%%%%%%%%%%%%%%%%%%%%%%%%%%%%%%%%
%% दिवस २: ध्यान घडत आहे का?
%%%%%%%%%%%%%%%%%%%%%%%%%%%%%%%%%%%%%%%%%%%%%%%%%%%%%%%%%%%%%%%%%%%%%%%%%%%%%%%%%%

%%%%%%%%%%%%%%%%%%%%%%%%%%%%%%%%%%%%%%%%%%%%%%%%%%%%%%%%%%%
\begin{frame}[fragile]\frametitle{दिवस २ --- पंखा उदाहरण}

	\begin{itemize}
	\item पंखा थांबवायला पंखेच थांबवत नाहीत --- स्विच बंद करतात
	\item मनावर थेट नियंत्रण शक्य नाही --- योग्य परिस्थिती निर्माण करायची
	\item विचार = अभौतिक; श्वास = त्यांचे भौतिक रूप
	\item मूर्तिपूजा उदाहरण: सूक्ष्म देवतेपर्यंत पोहोचण्यासाठी भौतिक मूर्तीची पूजा
	\item श्वासावर लक्ष केंद्रित केल्यास विचार आपोआप कमी होतात
	\item श्वास हळू करण्याचा सोपा मार्ग: मनात शांतपणे `ॐ' उच्चारा --- श्वासाला ॐ च्या लांबीचा वेळ द्या
	\end{itemize}

{\tiny (Ref: राकेश देशपांडे --- दिवस २)}

\end{frame}

%%%%%%%%%%%%%%%%%%%%%%%%%%%%%%%%%%%%%%%%%%%%%%%%%%%%%%%%%%%
\begin{frame}[fragile]\frametitle{दिवस २ --- स्थिरतेचा क्रम}

	\begin{itemize}
	\item निर्विचारता $\leftarrow$ हळू श्वास $\leftarrow$ स्थिर शरीर
	\item उलट क्रमाने काम करा: \textbf{शरीर स्थिर} $\rightarrow$ श्वास हळू $\rightarrow$ मन निर्विचार
	\item हस्तमुद्रेचा मुख्य हेतू: शरीर स्थिर ठेवणे
	\item ``स्थिरं सुखं आसनम्'' --- स्थिर आसन हेच खरे सुख देणारे
	\item अष्टांग योग (पतंजली): स्थिर आसन ध्यानासाठी आवश्यक
	\item हठयोग (स्वात्मराम, १४वे शतक): शारीरिक साधनांवर भर
	\end{itemize}

{\tiny (Ref: राकेश देशपांडे --- दिवस २)}

\end{frame}

%%%%%%%%%%%%%%%%%%%%%%%%%%%%%%%%%%%%%%%%%%%%%%%%%%%%%%%%%%%
\begin{frame}[fragile]\frametitle{दिवस २ --- शून्यावस्था + साधना}

	\begin{itemize}
	\item शून्यावस्था: ध्यानात एका क्षणासाठी श्वास थांबतो --- मन निर्विचार होते
	\item त्या क्षणाची जाणीव होते --- पण जाणीवही एक विचार असल्याने मन पुन्हा वळते
	\item सरावाने हा क्षण ५ सेकंदांपर्यंत टिकतो --- पुढे वाढतो
	\item झोप vs ध्यान: झोपेत श्वास चालू पण जाणीव नाही; ध्यानात स्वतः चालवतो
	\item स्वप्नात जाणीव असते की स्वप्न पाहत आहे --- ध्यानाच्या जवळची पण वेगळी अवस्था
	\end{itemize}

	\vspace{2mm}
	\textbf{साधना:} पाठीला आधार न घेता सरळ बसा, डोळे उघडे, २० मिनिटे स्थिर \\
	खाज सुटली तरी सहन करा; फक्त पापण्या हलवणे व लाळ गिळणे मान्य

{\tiny (Ref: राकेश देशपांडे --- दिवस २)}

\end{frame}

%%%%%%%%%%%%%%%%%%%%%%%%%%%%%%%%%%%%%%%%%%%%%%%%%%%%%%%%%%%%%%%%%%%%%%%%%%%%%%%%%%
%% दिवस ३: ध्यानातील अडथळे
%%%%%%%%%%%%%%%%%%%%%%%%%%%%%%%%%%%%%%%%%%%%%%%%%%%%%%%%%%%%%%%%%%%%%%%%%%%%%%%%%%

%%%%%%%%%%%%%%%%%%%%%%%%%%%%%%%%%%%%%%%%%%%%%%%%%%%%%%%%%%%
\begin{frame}[fragile]\frametitle{दिवस ३ --- विचार येणे स्वाभाविक}

	\begin{itemize}
	\item ध्यान करत असताना विचार येणे ही नैसर्गिक बाब आहे
	\item शरीर सुदृढ करायला जसा वेळ लागतो --- ध्यानालाही लागतो
	\item ध्यान सुरू केल्याच्या पहिल्याच दिवसापासून काही सकारात्मक परिणाम जाणवतात
	\item एका दिवशी मन अधिक शांत --- दुसऱ्या दिवशी विचारांचा पूर; हे स्वाभाविक आहे
	\item ध्यान फक्त बसण्याच्या वेळेपुरते नाही --- जागरूकता हीच जीवनशैली आहे
	\end{itemize}

{\tiny (Ref: राकेश देशपांडे --- दिवस ३)}

\end{frame}

%%%%%%%%%%%%%%%%%%%%%%%%%%%%%%%%%%%%%%%%%%%%%%%%%%%%%%%%%%%
\begin{frame}[fragile]\frametitle{दिवस ३ --- इतरांचे अनुभव + कंटाळा}

	\begin{itemize}
	\item इतरांना आलेल्या अनुभवाची अपेक्षा करणे चुकीचे --- प्रत्येकाचा अनुभव वेगळा
	\item प्रत्येकाने आपापल्या क्षमतेनुसार साधना करावी --- ध्यानासाठी कोणते सूत्र नाही
	\item कंटाळा येणे सामान्य --- ध्यान टाळण्याचा प्रयत्न होतो
	\item जर ३ महिन्यांनंतरही कंटाळाच येत असेल तर ती साधना कदाचित योग्य नाही
	\item तीच वेळ, तीच पद्धत पाळणे आवश्यक --- वारंवार पद्धत बदलणे चुकीचे
	\item सातत्य आणि निष्ठाच साधनेला परिपक्व करते
	\end{itemize}

{\tiny (Ref: राकेश देशपांडे --- दिवस ३)}

\end{frame}

%%%%%%%%%%%%%%%%%%%%%%%%%%%%%%%%%%%%%%%%%%%%%%%%%%%%%%%%%%%
\begin{frame}[fragile]\frametitle{दिवस ३ --- साधना: साक्षीभाव}

	\begin{itemize}
	\item दृश्य: भारत--पाकिस्तान क्रिकेट सामन्याच्या शेवटच्या चेंडूवर अचानक दिवे जातात
	\item त्या क्षणी प्रत्येकाचा श्वास रोखलेला --- एकाग्रतेने श्वास आपोआप मंदावतो
	\item ध्यानात अंतर्गत घटकांवर लक्ष केंद्रित केल्यास हेच घडते
	\item बाह्य परिस्थितीपेक्षा अंतर्गत निरीक्षण अधिक महत्त्वाचे
	\end{itemize}

	\vspace{2mm}
	\textbf{साधना:} डोळे बंद, स्थिर बसा, २० मिनिटे \\
	श्वास, हृदयाची धडधड, अनहद नाद यावर लक्ष --- साक्षीभाव विकसित करा

{\tiny (Ref: राकेश देशपांडे --- दिवस ३)}

\end{frame}

%%%%%%%%%%%%%%%%%%%%%%%%%%%%%%%%%%%%%%%%%%%%%%%%%%%%%%%%%%%%%%%%%%%%%%%%%%%%%%%%%%
%% दिवस ४: ध्यान मदत करत आहे का?
%%%%%%%%%%%%%%%%%%%%%%%%%%%%%%%%%%%%%%%%%%%%%%%%%%%%%%%%%%%%%%%%%%%%%%%%%%%%%%%%%%

%%%%%%%%%%%%%%%%%%%%%%%%%%%%%%%%%%%%%%%%%%%%%%%%%%%%%%%%%%%
\begin{frame}[fragile]\frametitle{दिवस ४ --- चिन्हे १ ते ३}

	\begin{itemize}
	\item \textbf{आंतरिक स्थिती बाह्य परिस्थितीवर अवलंबून राहात नाही}
		\begin{itemize}
		\item कोणावर प्रतिक्रिया द्यायची हे तुमच्या हातात असते
		\item बाह्य घटनांशी मन कसे जोडायचे हे तुम्ही ठरवू शकता
		\end{itemize}
	\item \textbf{स्वीकार (Acceptance)}: जे बदलता येत नाही ते स्वीकारणे; ऊर्जा नियंत्रणात असलेल्यावर केंद्रित करणे
	\item \textbf{अनासक्ती (Detachment)}: अतिरेकी जोडणे टाळणे
		\begin{itemize}
		\item आर्थिक नुकसान टाळण्यासाठी `स्टॉप लॉस' --- तसेच मानसिक हानी टाळणे
		\end{itemize}
	\end{itemize}

{\tiny (Ref: राकेश देशपांडे --- दिवस ४)}

\end{frame}

%%%%%%%%%%%%%%%%%%%%%%%%%%%%%%%%%%%%%%%%%%%%%%%%%%%%%%%%%%%
\begin{frame}[fragile]\frametitle{दिवस ४ --- चिन्हे ४ ते ७}

	\begin{itemize}
	\item \textbf{वर्तमान क्षणात राहण्याची क्षमता}: भूत-भविष्यात न अडकता जाणीवपूर्वक जगणे
	\item \textbf{सकारात्मकता वाढते}: दिवसाच्या २०\% विचार सकारात्मक झाले तर बाकीचेही त्याच दिशेने जातात; गोष्टी प्रत्यक्षात यायला लागतात
	\item \textbf{आसपासच्या लोकांमध्ये बदल}: नकारात्मक लोक दूर जातात; सकारात्मक येतात; तुमच्या ऊर्जेचा प्रभाव इतरांवर पडतो
	\item \textbf{भावना नियंत्रणात}: राग, भीती, दुःख निरीक्षणात आणता येतात; दडपण्याऐवजी स्वीकारून त्यावर काम करणे
	\end{itemize}

{\tiny (Ref: राकेश देशपांडे --- दिवस ४)}

\end{frame}

%%%%%%%%%%%%%%%%%%%%%%%%%%%%%%%%%%%%%%%%%%%%%%%%%%%%%%%%%%%
\begin{frame}[fragile]\frametitle{दिवस ४ --- सारांश}

	\begin{itemize}
	\item ध्यान एका दिवसात पूर्ण परिणाम देत नाही
	\item सातत्याने अभ्यास केल्यास जीवनाच्या प्रत्येक अंगावर सकारात्मक परिणाम घडतो
	\item मन अधिक स्थिर, जागरूक आणि शांत होत जाते
	\item जीवनाकडे पाहण्याचा दृष्टिकोनच बदलतो --- हीच खरी प्रगती आहे
	\end{itemize}

{\tiny (Ref: राकेश देशपांडे --- दिवस ४)}

\end{frame}

%%%%%%%%%%%%%%%%%%%%%%%%%%%%%%%%%%%%%%%%%%%%%%%%%%%%%%%%%%%%%%%%%%%%%%%%%%%%%%%%%%
%% दिवस ५: ध्यानाचे नियम व सूचना
%%%%%%%%%%%%%%%%%%%%%%%%%%%%%%%%%%%%%%%%%%%%%%%%%%%%%%%%%%%%%%%%%%%%%%%%%%%%%%%%%%

%%%%%%%%%%%%%%%%%%%%%%%%%%%%%%%%%%%%%%%%%%%%%%%%%%%%%%%%%%%
\begin{frame}[fragile]\frametitle{दिवस ५ --- तयारी आणि वेळ}

	\begin{itemize}
	\item ध्यानापूर्वी हलकी शारीरिक हालचाल + श्वसन व्यायाम (२--३ मिनिटे)
	\item दररोज एकाच वेळेला, दोन वेळा, प्रत्येकी सुमारे २० मिनिटे
	\item जेवणाच्या किमान २ तास आधी ध्यान करावे
	\item झोपण्यापूर्वी लगेच ध्यान टाळावे --- ध्यानाने मन ताजेतवाने होते
	\item संध्याकाळी शक्यतो ७ वाजण्यापूर्वी बसणे हितावह
	\item पोट रिकामे असणे आवश्यक; शरीराच्या सर्व नैसर्गिक गरजा आधी पूर्ण कराव्यात
	\end{itemize}

{\tiny (Ref: राकेश देशपांडे --- दिवस ५)}

\end{frame}

%%%%%%%%%%%%%%%%%%%%%%%%%%%%%%%%%%%%%%%%%%%%%%%%%%%%%%%%%%%
\begin{frame}[fragile]\frametitle{दिवस ५ --- आहार, नैतिकता, जीवनशैली}

	\begin{itemize}
	\item आहार सात्त्विक असावा --- स्वच्छ, पचनास सुलभ, अल्प मसाल्याचा
	\item `सात्त्विक' म्हणजे फक्त शाकाहारी नव्हे --- प्रकृतीला अनुरूप, ताज्या पदार्थांचा आहार
	\item पूर्वकृती: ध्यानापूर्वीचे विचार, कार्य, वातावरण --- यांचा थेट परिणाम ध्यानावर होतो
	\item उत्तरकृती: ध्यानानंतर मिळालेली शांती पुढील कृतींवर प्रभाव टाकते
	\item नैतिकतेची पायाभरणी आवश्यक --- असत्य नको, भ्रष्टाचार नको, समयपालन
	\item अनैतिक वर्तनाने अंतर्मनाला वेदना होतात --- एकाग्रतेला अडथळा निर्माण होतो
	\end{itemize}

{\tiny (Ref: राकेश देशपांडे --- दिवस ५)}

\end{frame}

%%%%%%%%%%%%%%%%%%%%%%%%%%%%%%%%%%%%%%%%%%%%%%%%%%%%%%%%%%%
\begin{frame}[fragile]\frametitle{दिवस ५ --- आशय निवड + सातत्य}

	\begin{itemize}
	\item दिवसभर ग्रहण केलेले --- श्रवण, वाचन, दर्शन --- सर्व मनावर नोंदवले जाते
	\item हिंसात्मक, नकारात्मक, खिन्न करणारा आशय टाळावा
	\item सोशल मीडियावरील आशयाची निवड सजगपणे करणे गरजेचे
	\item अस्थिर करणारे लोक, आठवणी, ठिकाणे, वस्तू यांपासून विवेकी माघार घेणे
	\item हे कमजोरीचे लक्षण नसून परिपक्वतेचे प्रतीक आहे
	\item तीच वेळ, तीच पद्धत --- दिवसातून दोन वेळा दृढ इच्छाशक्तीने साधना करणे हाच उपाय
	\end{itemize}

{\tiny (Ref: राकेश देशपांडे --- दिवस ५)}

\end{frame}

%%%%%%%%%%%%%%%%%%%%%%%%%%%%%%%%%%%%%%%%%%%%%%%%%%%%%%%%%%%%%%%%%%%%%%%%%%%%%%%%%%
%% दिवस ६: अवचेतन मनात कसे प्रवेश करावा?
%%%%%%%%%%%%%%%%%%%%%%%%%%%%%%%%%%%%%%%%%%%%%%%%%%%%%%%%%%%%%%%%%%%%%%%%%%%%%%%%%%

%%%%%%%%%%%%%%%%%%%%%%%%%%%%%%%%%%%%%%%%%%%%%%%%%%%%%%%%%%%
\begin{frame}[fragile]\frametitle{दिवस ६ --- संकल्प, अफर्मेशन, मॅनिफेस्टेशन}

	\begin{itemize}
	\item एखादी इच्छा अवचेतन मनात गेली की प्रक्रिया पूर्ण --- चांगली असो वा वाईट असो
	\item आध्यात्मिक साधना बाह्य मनाला मवाळ बनवते --- संकल्प सहज आत पोहोचतो
	\item इच्छा = \textbf{आकर्षणाचा नियम} (Law of Attraction)
	\item वारंवार दृढ करणे = \textbf{अफर्मेशन}
	\item प्रत्यक्षात येणे = \textbf{मॅनिफेस्टेशन}
	\item त्राटक साधना थेट अवचेतन मनाला बळकटी देते
	\end{itemize}

{\tiny (Ref: राकेश देशपांडे --- दिवस ६)}

\end{frame}

%%%%%%%%%%%%%%%%%%%%%%%%%%%%%%%%%%%%%%%%%%%%%%%%%%%%%%%%%%%
\begin{frame}[fragile]\frametitle{दिवस ६ --- युगे आणि जीवनाचा उद्देश}

	\begin{itemize}
	\item सत्य युग: मत्स्य, कूर्म, वराह, नरसिंह --- सृष्टीचे संरक्षण
	\item त्रेता युग: वामन, परशुराम, राम --- न्यायाची पुनर्स्थापना
	\item द्वापर युग: कृष्ण (प्रेम, भक्ती, योग), बुद्ध (करुणा, ध्यान)
	\item कलियुग: भविष्यात कल्की --- अधर्माचा नाश
	\item प्रत्येक युगात चांगले-वाईट दोन्ही असते --- विवेकाने निवडक नैतिकता
	\item जीवनाचा सर्वात महत्त्वाचा उद्देश: \textbf{स्वतःला जाणून घेणे}
	\item आपल्या मनावर मानसिक सवयी, विचार, भावना यांचे थर चढतात --- ते विरळणे आवश्यक
	\end{itemize}

{\tiny (Ref: राकेश देशपांडे --- दिवस ६)}

\end{frame}

%%%%%%%%%%%%%%%%%%%%%%%%%%%%%%%%%%%%%%%%%%%%%%%%%%%%%%%%%%%
\begin{frame}[fragile]\frametitle{दिवस ६ --- त्राटक: बिंदु त्राटक}

	\begin{itemize}
	\item ८$\times$८ इंच पांढरा कागद, २ इंच व्यासाचा काळा बिंदू, ३ फूट अंतर
	\item डोळे न मिचकावता ध्यान --- सुरुवातीस ५ मिनिटे; हळूहळू २० मिनिटांपर्यंत
	\item काळा बिंदू $\rightarrow$ प्रकाशमान बिंदू $\rightarrow$ नाहीसा होतो = अवचेतन मनात प्रवेश
	\item साधना संपल्यानंतर डोळे स्वच्छ पाण्याने धुणे; संध्याकाळी करणे योग्य
	\item \textbf{दुसरा टप्पा}: पहिल्या महिन्यात १ इंच बिंदू, पुढे अर्धा, तिसऱ्या महिन्यात अतिशय लहान
	\end{itemize}

{\tiny (Ref: राकेश देशपांडे --- दिवस ६)}

\end{frame}

%%%%%%%%%%%%%%%%%%%%%%%%%%%%%%%%%%%%%%%%%%%%%%%%%%%%%%%%%%%
\begin{frame}[fragile]\frametitle{दिवस ६ --- प्रतिबिंब त्राटक व ज्योती त्राटक}

	\begin{itemize}
	\item \textbf{प्रतिबिंब त्राटक:} ८$\times$८ इंच आरसा, ३ फूट अंतर
		\begin{itemize}
		\item दोन्ही भुवयांच्या मधोमध लक्ष केंद्रित करावे
		\item चेहरा हळूहळू विरघळल्यासारखा दिसतो = अंतर्मनाशी नाते जुळले
		\end{itemize}
	\item \textbf{ज्योती त्राटक:} मेणबत्ती एकमेव प्रकाश; इतर सर्व प्रकाश बंद
		\begin{itemize}
		\item सुरुवातीस ५ मिनिटे; हळूहळू २० मिनिटांपर्यंत
		\item दिवसातून दोन वेळा सराव
		\end{itemize}
	\item या सर्व साधना संयमाने केल्यास अवचेतन मनाशी असलेले नाते दृढ होते
	\end{itemize}

{\tiny (Ref: राकेश देशपांडे --- दिवस ६)}

\end{frame}

%%%%%%%%%%%%%%%%%%%%%%%%%%%%%%%%%%%%%%%%%%%%%%%%%%%%%%%%%%%%%%%%%%%%%%%%%%%%%%%%%%
%% दिवस ७: या जीवनाचा खरा शोध
%%%%%%%%%%%%%%%%%%%%%%%%%%%%%%%%%%%%%%%%%%%%%%%%%%%%%%%%%%%%%%%%%%%%%%%%%%%%%%%%%%

%%%%%%%%%%%%%%%%%%%%%%%%%%%%%%%%%%%%%%%%%%%%%%%%%%%%%%%%%%%
\begin{frame}[fragile]\frametitle{दिवस ७ --- ध्यान आणि आत्मस्वरूप}

	\begin{itemize}
	\item जीवनाचा मुख्य उद्देश: स्वतःच्या खऱ्या स्वरूपाचा शोध घेणे
	\item ध्यानाची प्राथमिक व्याख्या: विचारशून्यता प्राप्त करणे --- पण ध्यान येथेच थांबत नाही
	\item विचार थांबल्यानंतर आत्म्याच्या गूढ शोधाची आध्यात्मिक यात्रा सुरू होते
	\item या अवस्थेला आध्यात्मिक ध्यान म्हणतात --- मनाच्या मर्यादा ओलांडणे
	\item शून्यावस्थेत आत्मस्वरूप प्रकट होते --- आत्मा मनाच्या पलीकडे असतो
	\end{itemize}

{\tiny (Ref: राकेश देशपांडे --- दिवस ७)}

\end{frame}

%%%%%%%%%%%%%%%%%%%%%%%%%%%%%%%%%%%%%%%%%%%%%%%%%%%%%%%%%%%
\begin{frame}[fragile]\frametitle{दिवस ७ --- मन, वेळ आणि नाद}

	\begin{itemize}
	\item मन हे वेळ आणि अवकाश निर्माण करणारे --- मन थांबते तिथे वेळ-अवकाशही नाही
	\item नाद/ध्वनी/कंपन: वेळ आणि अवकाशाच्या बंधनांपलीकडे असते
	\item अवकाशाच्या पलीकडे = कोणतेही ठिकाण नाही, सीमा नाही, आकार नाही --- सर्वत्र अस्तित्व
	\item वेळेच्या पलीकडे = कारण-परिणाम सिद्धांत लागत नाही --- सुरुवात नाही, शेवट नाही
	\item या अवस्थेला \textbf{स्थलकालातीत} (Transcending Time \& Space) म्हणतात
	\end{itemize}

{\tiny (Ref: राकेश देशपांडे --- दिवस ७)}

\end{frame}

%%%%%%%%%%%%%%%%%%%%%%%%%%%%%%%%%%%%%%%%%%%%%%%%%%%%%%%%%%%
\begin{frame}[fragile]\frametitle{दिवस ७ --- जीवनाचा खरा शोध}

	\begin{itemize}
	\item मनाच्या पलीकडे असलेला नाद प्रत्येक क्षणी, प्रत्येक ठिकाणी असतो
	\item जेव्हा विचारांपासून मुक्त होतो तेव्हा आत्म्याची अनुभूती संभवते
	\item ध्यान: केवळ मन शांत करणे नाही --- अस्तित्वाच्या परम सत्याशी जोडणारा सेतू
	\item बाह्य घटनांपासून मन मागे घेत आत्म्याच्या उपस्थितीत स्थिर राहण्याची क्षमता विकसित होते
	\item स्वतःला जाणून घेणेच जीवनाचे अंतिम ध्येय --- ही एक अद्भुत यात्रा आहे
	\end{itemize}

{\tiny (Ref: राकेश देशपांडे --- दिवस ७)}

\end{frame}
